\chapter{Pauli Spinor Lift of the Projective Operator Algebra}
\label{chap:pauli-spinor-lift}

\status{Exact CP$^1$/SU(2) double-cover algebra proved. No new xi symmetry, zero-location theorem, or RH claim is made.}

Chapter~\ref{chap:euler-riemann-negative-inversion} established the projective
operator group
\[
 R(u)=\frac1u,
 \qquad
 E(u)=-u,
 \qquad
 N(u)=-\frac1u,
\]
with
\[
 R^2=E^2=N^2=I,
 \qquad
 RE=ER=N,
\]
and hence
\[
 \{I,R,E,N\}\cong V_4.
\]
This chapter lifts that projective algebra to determinant-one unitary Pauli
matrices and identifies the exact spinorial group lying above it.

\section{Pauli representatives of the projective maps}

Use the Pauli matrices
\[
 \sigma_x=
 \begin{pmatrix}0&1\\1&0\end{pmatrix},
 \qquad
 \sigma_y=
 \begin{pmatrix}0&-i\\i&0\end{pmatrix},
 \qquad
 \sigma_z=
 \begin{pmatrix}1&0\\0&-1\end{pmatrix}.
\]
A matrix
\[
 M=\begin{pmatrix}a&b\\c&d\end{pmatrix}
\]
acts on the affine chart of \(\mathbb{CP}^1\) through the Möbius map
\[
 u\longmapsto\frac{au+b}{cu+d}.
\]
Since multiplying \(M\) by a non-zero scalar does not change its projective
action, the G012 maps admit representatives
\[
 R\sim\sigma_x,
 \qquad
 E\sim\sigma_z,
 \qquad
 N\sim-i\sigma_y.
\]

\section{Determinant-one unitary lifts}

Choose
\begin{equation}
 \boxed{\widetilde R=i\sigma_x},
 \qquad
 \boxed{\widetilde E=i\sigma_z},
 \qquad
 \boxed{\widetilde N=i\sigma_y}.
 \label{eq:g013-lifts}
\end{equation}
Each matrix satisfies
\[
 \widetilde A^\dagger\widetilde A=I,
 \qquad
 \det\widetilde A=1,
\]
so all three belong to \(SU(2)\).  Their projective actions are exactly
\[
 \widetilde R:u\mapsto\frac1u,
 \qquad
 \widetilde E:u\mapsto-u,
 \qquad
 \widetilde N:u\mapsto-\frac1u.
\]
Thus they are genuine spinor lifts of the projective maps in G012.

\section{Quaternion multiplication above projective commutativity}

The Pauli multiplication law gives
\begin{align}
 \widetilde R\widetilde E&=\widetilde N,
 &\widetilde E\widetilde R&=-\widetilde N,\\
 \widetilde E\widetilde N&=\widetilde R,
 &\widetilde N\widetilde E&=-\widetilde R,\\
 \widetilde N\widetilde R&=\widetilde E,
 &\widetilde R\widetilde N&=-\widetilde E.
\end{align}
Moreover,
\begin{equation}
 \boxed{
 \widetilde R^2
 =\widetilde E^2
 =\widetilde N^2
 =-I.
 }
 \label{eq:g013-square-minus-I}
\end{equation}
Hence every non-central generator has order four in the lifted group.

\begin{theorem}[Quaternion spinor lift]
The subgroup of \(SU(2)\) generated by \(\widetilde R\) and
\(\widetilde E\) is
\[
 \boxed{
 Q_8=\{\pm I,\pm i\sigma_x,\pm i\sigma_y,\pm i\sigma_z\}.
 }
\]
\end{theorem}

\begin{proof}
Equation~\eqref{eq:g013-square-minus-I} supplies the central element \(-I\),
while the displayed multiplication identities give closure of the eight
listed matrices.  They are pairwise distinct as matrices and satisfy the
standard quaternion multiplication table.  Therefore the generated subgroup
is \(Q_8\).
\end{proof}

The apparent conflict with the commutativity proved in G012 is resolved by
projectivization.  Matrix products satisfy
\[
 \widetilde R\widetilde E
 =-\widetilde E\widetilde R,
\]
but \(M\) and \(-M\) induce the same Möbius transformation.  Thus the central
spinor sign disappears on \(\mathbb{CP}^1\).

\section{The exact quotient to the G012 Klein group}

The center of the quaternion group is
\[
 Z(Q_8)=\{\pm I\}.
\]
Projectivization has this center as its kernel, and the four resulting classes
are exactly \(I,R,E,N\).  Consequently
\begin{equation}
 \boxed{Q_8/\{\pm I\}\cong V_4.}
 \label{eq:g013-q8-v4}
\end{equation}
Equivalently, the operator algebra sits in the short exact sequence
\begin{equation}
 \boxed{
 1\longrightarrow\{\pm I\}
 \longrightarrow Q_8
 \longrightarrow V_4
 \longrightarrow1.
 }
 \label{eq:g013-short-exact}
\end{equation}
This identifies precisely which information is lost when passing from the
spinor lift to its projective action: the central sign \(-I\).

\section{Spinorial order doubling}

A standard Pauli rotation through angle \(\pi\) about axis \(j\) is
\[
 U_j(\pi)
 =\exp\!\left(-\frac{i\pi}{2}\sigma_j\right)
 =-i\sigma_j.
\]
It differs from the chosen lift \(+i\sigma_j\) only by \(-I\), and therefore
induces the same CP$^1$ transformation.  Nevertheless,
\[
 U_j(\pi)^2=-I,
 \qquad
 U_j(\pi)^4=I.
\]
Thus the projective operation has order two while its spinor lift has order
four.  Two projective half-turns produce the central spinor phase \(-I\); four
half-turns are required to return the lifted spinor to the identity matrix.

\begin{corollary}[Projective order two, spinor order four]
Each non-trivial involution in the G012 Klein group has a Pauli spinor lift of
order four, and its square is the central element \(-I\).
\end{corollary}

\section{Fixed pairs as the three Bloch axes}

Use the homogeneous spinor chart \((u,1)^T\).  Its Bloch vector is
\[
 \mathbf b(u)
 =\frac{1}{1+|u|^2}
 \left(
 2\Re u,
 -2\Im u,
 |u|^2-1
 \right),
\]
with \(u=\infty\) mapped to \((0,0,1)\).

The three projective Pauli operators have fixed pairs
\[
 R:\quad u=\pm1,
\]
\[
 E:\quad u=0,\infty,
\]
\[
 N:\quad u=\pm i.
\]
Their Bloch images are the three antipodal coordinate-axis pairs:
\[
 R\leftrightarrow\pm x,
 \qquad
 E\leftrightarrow\pm z,
 \qquad
 N\leftrightarrow\pm y,
\]
where the sign orientation of the \(y\)-axis follows the chosen homogeneous
chart convention.

In particular, the G012 negative-inversion fixed pair
\[
 u=\pm i
\]
is precisely the Pauli-\(y\) eigenline pair.  Under
\(s=u/(1+u)\), these two projective points remain
\[
 \boxed{s=\frac12\pm\frac{i}{2}},
\]
as established in Chapter~37.

\section{Proof firewall}

Proved in SOH-G013:
\begin{itemize}
 \item the Pauli representatives and the determinant-one unitary lifts in
       \cref{eq:g013-lifts};
 \item quaternion multiplication and \(Q_8\) closure;
 \item the order-four relation \(\widetilde A^2=-I\);
 \item the exact projective quotient
       \(Q_8/\{\pm I\}\cong V_4\);
 \item the short exact sequence \cref{eq:g013-short-exact};
 \item the central-sign origin of projective commutativity;
 \item the equivalence, up to \(-I\), with the standard Pauli \(\pi\)-rotation
       lifts;
 \item the three antipodal Bloch-axis fixed pairs.
\end{itemize}

Not proved in SOH-G013:
\begin{itemize}
 \item any new functional equation for \(\xi\);
 \item invariance of the xi zero set under \(E\) or \(N\);
 \item that a non-trivial xi zero must be a Pauli fixed line;
 \item PF\(_\infty\);
 \item SOH-G003 real-rootedness;
 \item the Riemann Hypothesis.
\end{itemize}
